\documentclass[11pt]{article}
\usepackage[a4paper,margin=1in]{geometry}
\usepackage{amsmath,amssymb,mathtools,bm}
\usepackage{enumitem,booktabs,array}
\usepackage{tcolorbox}
\usepackage{hyperref}
\usepackage[protrusion=true,expansion=false]{microtype}
\hypersetup{colorlinks=true,linkcolor=blue,urlcolor=blue,citecolor=blue}
\setlist[itemize]{topsep=2pt,itemsep=2pt}
\setlist[enumerate]{topsep=2pt,itemsep=2pt}
\newtcolorbox{keybox}{colback=blue!3!white,colframe=blue!40!black,boxrule=0.4pt,arc=1mm}
\newtcolorbox{alertbox}{colback=red!3!white,colframe=red!40!black,boxrule=0.4pt,arc=1mm}
\newtcolorbox{answerbox}{colback=green!3!white,colframe=green!40!black,boxrule=0.4pt,arc=1mm}
\newtcolorbox{drillbox}{colback=yellow!5!white,colframe=orange!60!black,boxrule=0.4pt,arc=1mm}
\newcommand{\EV}{\mathbb{E}}
\newcommand{\Var}{\mathrm{Var}}
\newcommand{\Cov}{\mathrm{Cov}}
\newcommand{\blank}[1]{\underline{\hspace{#1}}}

\title{E06-04: Kepler, Naiker \& Stewart (2023, JF) \\
\large ``Stealth Acquisitions and Product Market Competition'' --- Active Recall Workbook}
\author{Empirical Corporate Finance II --- Exam Prep}
\date{\today}
\begin{document}\maketitle

\begin{keybox}
\textbf{Paper in one sentence.} KNS show that a meaningful share of M\&A falls below disclosure thresholds (``stealth acquisitions'') and therefore escapes antitrust scrutiny, and that these small deals, when aggregated in product-market networks, measurably reduce competition and increase markups --- evidence that the Hart-Scott-Rodino (HSR) reporting threshold creates an enforcement blind spot.

\textbf{Placement.} Recent paper on M\&A and competition policy; connects to Azar-Schmalz-Tecu common-ownership literature and antitrust enforcement debates. Important for applied antitrust research.
\end{keybox}

\section*{Round 1: Complete Walk-Through}

\subsection*{1.1 Institutional background}
\begin{itemize}
\item Hart-Scott-Rodino Act (1976) requires pre-merger notification for deals above a size threshold (currently $\sim$\$100 million).
\item Below-threshold deals are not centrally reported and face minimal antitrust scrutiny.
\item ``Stealth acquisitions'' = cumulative sub-threshold acquisitions by a firm over time, which can add up to significant market concentration.
\end{itemize}

\subsection*{1.2 Data}
\begin{itemize}
\item SDC, PitchBook, and proprietary data on sub-threshold deals, 2000--2020.
\item TNIC industry classification to define product-market neighborhoods.
\item Firm-level markups (De Loecker-Eeckhout method), sales concentration, R\&D, entry.
\end{itemize}

\subsection*{1.3 Empirical design}
\begin{itemize}
\item Industry-level panel: within-industry stealth M\&A intensity as treatment.
\item Compare industries with high vs.\ low stealth activity on markups, concentration, entry.
\item Firm-level: acquirers with repeated stealth deals vs.\ matched controls.
\item Event-study around HSR threshold changes as identification.
\end{itemize}

\subsection*{1.4 Headline results}
\begin{itemize}
\item Stealth acquisitions represent $\sim$30\% of M\&A volume by count, though much smaller by value.
\item Industries with concentrated stealth acquirers see markups rise by 2--4 pp over 5 years.
\item Product-market competitors report lower growth; entry falls by $\sim$15\%.
\item HSR threshold increases (raising the reporting bar) predict spikes in below-threshold acquisitions just under the new cutoff.
\end{itemize}

\subsection*{1.5 Mechanism}
\begin{itemize}
\item Serial small acquisitions create market power invisible to antitrust review.
\item Startups and potential competitors are bought before they can scale.
\item Product-market network effects amplify individual-deal anticompetitive impact.
\end{itemize}

\subsection*{1.6 Implications}
\begin{itemize}
\item Antitrust enforcement should consider cumulative M\&A activity, not just individual-deal size.
\item Calls for sub-threshold reporting or industry-concentration triggers.
\item Evidence relevant to debates over Big Tech ``killer acquisitions'' (Cunningham-Ederer-Ma 2021).
\end{itemize}

\subsection*{1.7 Caveats}
\begin{alertbox}
\begin{itemize}
\item Sub-threshold deal data is incomplete; measurement error biases estimates.
\item Markup measurement (De Loecker-Eeckhout) is itself contested.
\item Identification relies on threshold variation that may coincide with other regulatory changes.
\end{itemize}
\end{alertbox}

\section*{Round 2: Level 1 Fill-in}
\begin{enumerate}
\item HSR reporting threshold: $\sim\$$\blank{1cm} million.
\item Stealth deals' share of M\&A: $\sim$\blank{1cm}\% by count.
\item Markup effect: $+$\blank{1cm}--\blank{1cm} pp over 5 years.
\item Entry effect: \blank{1cm}\% decline.
\item Related concept: \blank{2cm} acquisitions (Cunningham-Ederer-Ma).
\end{enumerate}

\section*{Round 3: Level 2 Prompts}
\begin{enumerate}
\item Describe the HSR threshold and stealth-M\&A definition.
\item Report main findings on markups, entry, and concentration.
\item Explain the mechanism and use of TNIC networks.
\item Discuss implications for antitrust policy.
\item How does KNS relate to killer-acquisitions research?
\end{enumerate}

\section*{Round 4: Minimal Scaffolding}
From a blank page: (i) institutional setting, (ii) data, (iii) findings, (iv) mechanism, (v) policy.

\section*{Round 5: Blank Page}
Essay: stealth acquisitions, market power, and antitrust enforcement.

\vspace{6pt}
\begin{drillbox}
\textbf{Speed Drill.} (1) Threshold. (2) Stealth share. (3) Markup effect. (4) Entry effect. (5) Policy implication.
\end{drillbox}

\section*{Quick Reference \& Answer Key}
\begin{answerbox}
\begin{itemize}
\item \textbf{Setting.} HSR threshold exempts small deals.
\item \textbf{Share.} $\sim$30\% of M\&A count sub-threshold.
\item \textbf{Effect.} Markups $+2$--$4$ pp; entry $-15\%$.
\item \textbf{Policy.} Enforce on cumulative activity.
\end{itemize}
\end{answerbox}

\begin{keybox}
\textbf{30-second exam pitch.} Kepler, Naiker \& Stewart (2023) document ``stealth acquisitions'' --- sub-Hart-Scott-Rodino-threshold deals that escape antitrust review. Using SDC, PitchBook, and proprietary deal data 2000--2020 combined with TNIC product-market networks and De Loecker-Eeckhout markups, they find stealth deals account for $\sim$30\% of M\&A by count. In industries where stealth acquirers concentrate, markups rise 2--4 pp over 5 years, competitor growth falls, and entry drops $\sim$15\%. HSR threshold increases predict spikes in deals sized just below the new cutoff. Mechanism: serial small acquisitions of potential competitors create market power invisible to individual-deal review. The paper motivates antitrust reforms that consider cumulative M\&A activity (not just single-deal size) and directly informs ``killer acquisitions'' debates, including Big Tech's serial startup buyouts. It is a canonical recent contribution at the intersection of M\&A and competition policy.
\end{keybox}

\end{document}
